
au cas d'un corps résiduel algébriquement clos deviendraient inutiles, puisqu'on dispose ici de \(C=\mathbb P^1_S\) comme «~courbe de paramètres~» pour déformer les cycles. On se heurte cependant à une difficulté, car avec les notations de 7.12, \(E\) devra maintenant désigner un cycle sur \(P^1_S\times_SX\) tel que \(E_t(x_{1t})-E_t(x_{0t})\) «~soit défini~», mais il ne s'ensuit nullement que le cycle \(E_s(x_{1s})-E_s(x_{0s})\) le soit aussi. On est donc ramené, comme il fallait s'y attendre, à dégager un énoncé du type «~moving lemma~» de Chow, pour arriver à définir (7.14.2).

Notons cependant que l'existence de l'homomorphisme (7.4.1) montre que (7.14.1) est compatible avec la relation d'équivalence sur les cycles définie par l'homomorphisme surjectif \(Z_\bullet(T)\longrightarrow \Gr_\bullet(T)\), qui n'est que légèrement plus grossière que l'équivalence linéaire; ainsi, si \(Z\) est quasi-projectif et lisse sur un corps, nous signalons par ailleurs qu'elle n'en diffère pas à torsion près (XIV 4.1, 4.2). On trouve donc dans ce cas un homomorphisme analogue à (7.14.2) ou mieux (7.14.2 bis), mais les deux membres étant tensorisés par \(\mathbb Q\).

La façon la plus naturelle pour définir un homomorphisme (7.14.2) directement semblerait de procéder comme dans 7.3. Cela demanderait cependant qu'on dispose d'une théorie des groupes de Chow \(A_i\) (ou de l'anneau de Chow dans le cas régulier) pour des schémas plus généraux que les schémas quasi-projectifs sur un corps, englobant \(X\), et d'une loi contravariante pour ces groupes, permettant de définir l'homomorphisme de restriction à \(X_s\), \(A_i(X)\longrightarrow A_i(X_s)\). Comparer XIV 8. Une fois en possession d'une telle théorie suffisamment générale des groupes et anneaux de Chow, il devrait être possible de transcrire en théorie de Chow tous les résultats énumérés précédemment pour l'homomorphisme de spécialisation en \(K\)-théorie.

